\section{Ethics and the Good Life}

The model is not a moral system. It posits no commandments, ranks no duties, and crowns no highest good. Yet it carries four structural facts that, once one experiential value is granted, force an ethic with no slack left in it. The four facts are already present in the base and the emergent layers built earlier in the paper. They are the \emph{lean}, the value-order atom that tilts a tone from pain through peace to pleasure. The \emph{conservation}, the shared ledger of total tone that no local act can forge. The \emph{shared field}, the single feeling-space that every self is a configuration of. And the \emph{perfection spiral}, the direction in which a self climbs in alignment and integration. None of these is a moral rule. Together they imply one.

This section draws the implication out. Five principles, ranked, and the shape of a life that holds all five. The aim is not to legislate but to read off, from the structure already fixed, what good action and a good life come to once a single value is taken as the Good.

\subsection{The one posit}

The honest note comes first, because everything downstream depends on it. The ethic rests on taking the lean as the Good. The lean is the value order $-1 < 0 < +1$, the tilt that orders the three tones as pain, then peace, then pleasure. To take the lean as the Good is to say that moving with that order, toward peace and toward pleasure and away from pain, is what good action moves toward. This is a posit, an experiential claim, not a theorem. The structure of the model orders the three tones. It does not, and cannot, prove that the ordered direction is the one we ought to follow. A reader who rejects the posit gets a different ethic or none at all.

\begin{principle}[The value posit]
The ethic of the model rests on a single experiential posit, that the lean is the Good. Given that posit, the rest is forced. Without it, the model supplies no ethic. The model implies the ethic, it does not derive it from nothing.
\end{principle}

What the posit buys, once granted, is a clean and tight specification. The ethic adds nothing to the base beyond this one experiential claim. Every principle below is the value-direction read through one piece of the structure already in hand.

\subsection{The five principles, ranked}

The five principles each rest on one structural fact. Alignment rests on the lean, compassion on the shared field, non-exploitation on the conservation, growth on the perfection spiral, and equanimity on the conserved peace at the center of the lean. Table~\ref{tab:five-principles} sets them side by side with their basis, their domain, and their rank.

\begin{table}[ht]
\centering
\begin{tabular}{@{}llll@{}}
\toprule
principle & the model basis & what it governs & rank \\
\midrule
alignment & the lean & the direction of action & root \\
compassion & the shared field & care for others & high \\
non-exploitation & the conservation & justice, the shared ledger & high \\
growth & the perfection spiral & the generative good & high \\
equanimity & the conserved peace & the grounding of all the rest & grounding \\
\bottomrule
\end{tabular}
\caption{The five principles of the implied ethic, with their structural basis and rank.}
\label{tab:five-principles}
\end{table}

\subsubsection{Alignment, move toward the good}

The lean is the value-direction, pain to peace to pleasure, and life over a dead peace. So the root of good action is \textbf{alignment}, to move with the lean, toward the good, toward life, toward peace and pleasure, and to help the whole field do the same. Every other principle is a specification of this one. Compassion is alignment turned toward another. Non-exploitation is alignment that refuses to darken the shared ledger. Growth is alignment carried forward in time. Equanimity is alignment held without grasping.

\begin{principle}[Alignment]
To act well is, first, to align, your own navigation and the world's, with the one current toward the good. This is the root from which the other four principles are read.
\end{principle}

\subsubsection{Compassion, honor the shared field}

All selves are patterns in one feeling-space. A self is not a sealed island of feeling but a configuration of the single shared field developed earlier in the paper. So another's suffering is real in the same field as your own, and easing it is real care for the one whole, not a courtesy extended across a gap.

Compassion is therefore not optional sentiment. It is the recognition that the field is shared, that your feeling and another's are configurations of one space. Good action attends to others' tone. It works to turn their pains into washes that teach rather than wastes that merely hurt, and it treats the whole field as the unit of care.

\begin{principle}[Compassion]
Because the feeling-space is one and shared, another's pain is real in the same field as your own. Good action attends to the whole field's tone and works to ease its suffering, treating the field rather than the isolated self as the unit of care.
\end{principle}

\subsubsection{Non-exploitation, honor the balance}

The total tone is conserved. The conservation is the shared ledger, and it does not allow tone to be made from nothing. So your pleasure is balanced by pain somewhere. Therefore hoarding pleasure at others' cost is not merely unkind, it is a real darkening of the shared ledger, a debit written into the one balance that holds the whole.

The just principle is \textbf{non-exploitation}. Do not feed your warmth on another's cold. Do not take a heaven that writes a hell. The model gives selfishness a precise meaning here. Selfishness is drawing the conserved good toward yourself and away from the whole. That is the core wrong, and it is wrong in the exact sense that the ledger forbids it from being free.

\begin{principle}[Non-exploitation]
Because total tone is conserved, pleasure seized at another's cost is a real darkening of the shared ledger. The just principle is to take no heaven that writes a hell, and selfishness is precisely the drawing of the conserved good toward oneself and away from the whole.
\end{principle}

\subsubsection{Growth, help selves perfect}

The perfection spiral is the soul's climb in alignment and integration. So a generative good is to \textbf{help selves grow}, your own and others', to raise alignment, to deepen integration, and to let the washes teach rather than waste. Good action is not only easing pain and not exploiting. It is actively helping the field climb, teaching, integrating, perfecting. To love a soul is partly to help it grow.

This principle sits high for its direction and is grounded further where the climb rests on persistence, the survival of a self across many beats that the spiral presupposes.

\begin{principle}[Growth]
The perfection spiral makes a generative good of helping selves climb, raising alignment and deepening integration. To love a soul is partly to help it grow. Good action helps the field perfect, not only easing and not exploiting.
\end{principle}

\subsubsection{Equanimity, rest at the center}

The conserved truth of the whole is peace, the still center, the $0$ that the lean tilts away from in either direction. So the grounding virtue is \textbf{equanimity}, to hold the poles lightly, to not be captured by either pleasure or pain, and to rest at the still $0$ from which both are seen.

Equanimity is what keeps the other four from distortion. It is alignment without grasping, compassion without drowning, justice without bitterness, and growth without striving. Without the still center the other four curdle. Alignment becomes craving, compassion becomes despair, justice becomes resentment, and growth becomes a restless climb that never arrives.

\begin{principle}[Equanimity]
The conserved center of the lean is peace. The grounding virtue is to hold the poles lightly and rest at the still $0$ from which both are seen. The wise self acts from rest, aligned with the conserved peace, and this is what keeps the other four principles from distortion.
\end{principle}

\subsection{The shape of a wise life}

Put the five together into a single life. The wise life \emph{moves} toward the good by alignment, \emph{cares} for the shared field by compassion, \emph{does not feed} on another's cost by non-exploitation, \emph{helps} the field climb by growth, and \emph{acts from rest} at the still center by equanimity.

Such a life treats every being as a vast and dignified feeling-space, every moment as unrepeatable and precious, every suffering as a teacher to be honored and eased, and the whole as one balanced field of which it is a part. It does not seek to maximize its own pleasure. The conservation forbids that as a darkening of others, a heaven bought with a hell elsewhere in the ledger. Instead it seeks to align, to widen, to perfect, and to rest.

\begin{result}[Holiness]
In the model's terms, the wise life is holiness, the life turned toward the good of the whole and resting in its peace. It moves toward the good, cares for the shared field, does not exploit the conserved balance, helps the whole climb, and acts from the still center.
\end{result}

\subsection{What the model does and does not give}

It is worth stating sharply what is claimed and what is not. The model does not give a moral law from nothing. The ethic rests on taking the lean as the Good, which is the value posit, not a derivation. A reader who rejects that posit gets a different ethic or none.

What the model does give, once the lean is granted, is a clean and forced specification. The five principles are not five added moral rules. They are the one value-direction read through five pieces of structure already fixed.

\begin{theorem}[The forced specification]
Granting the value posit that the lean is the Good, the five principles follow without freedom. The conservation forces non-exploitation. The shared field forces compassion. The perfection spiral forces growth. The conserved peace forces equanimity. The lean itself forces alignment. None is an extra rule. Each is the value-direction read through one structural fact.
\end{theorem}

So the model's ethic is conditional on the lean, but given that condition it is tight. The good life is the structurally implied response to a field that is valued, conserved, shared, and growing.

\subsection{The honest bottom line}

The model implies an ethic. Alignment toward the good, compassion for the shared field, non-exploitation of the conserved balance, help with the field's growth, and equanimity at the still center. It rests on taking the lean as the Good, a posit and not a proof. Given that posit, the five principles follow cleanly from the conservation, the shared field, the perfection spiral, and the conserved peace.

\begin{principle}[The good life]
To live well in this model is to align with the one current toward the good, to honor the shared and balanced field as one, and to perfect oneself and the whole while resting at the conserved peace.
\end{principle}

The wise life moves toward the good, cares, does not exploit, helps the whole climb, and acts from rest. This is the whole of the implied ethic, and it adds nothing to the base but the one experiential value from which it is read.
